\documentclass[a4paper,11pt]{jsarticle}
\usepackage[margin=1cm]{geometry}
\usepackage{amsmath, amssymb}
\usepackage{multicol}

% ページ番号を消す設定
\pagestyle{empty}

% ページを真ん中で分ける線の設定
\setlength{\columnsep}{1cm}
\setlength{\columnseprule}{0.4pt}

\newcommand{\testheader}{
\begin{center}
{\large 代数概論・同演習 中間テスト(11/16)}\\[1em]
\underline{\hspace{1cm}}年\underline{\hspace{1cm}}組\underline{\hspace{1cm}}番\hspace{0.5cm}名前\underline{\hspace{6cm}}
\end{center}
}

\begin{document}

% ==========================================
% 1枚目 表 (1ページ目: IMG_6467)
% ==========================================
\testheader

\begin{multicols*}{2}
% 領域1
\noindent
\textbf{問1.}\quad 次の言葉を説明せよ．\\[0.5em]
(1)[3点] 多項式の次数
\vspace{4cm}

\noindent
(2)[3点] $\operatorname{ord}_m(\bar{r})$, ただし $m \in \mathbb{N}$ で $\bar{r} \in (\mathbb{Z}/m\mathbb{Z})^{\times}$ とする．
\vspace{4cm}

\noindent
(3)[3点] フェルマーの小定理
\vspace{4cm}

\noindent
(4)[3点] 既約剰余類
\vspace{4cm}

\noindent
(5)[3点] 環 $R$ の既約元
\vspace{4cm}

\vfill\null
\columnbreak

% 領域2
\noindent
\textbf{問2.}[6点] $a, b, s \in \mathbb{Z}$, $m \in \mathbb{N}$ とする．$sa \equiv sb \pmod m$ で $\gcd(s, m) = 1$ ならば $a \equiv b \pmod m$ となることを示せ．
\vspace{11cm}

\noindent
\textbf{問3.}[6点] $\bar{r} \in (\mathbb{Z}/m\mathbb{Z})^{\times}$ に対し $l = \operatorname{ord}_m(\bar{r})$ とする．このとき，$\bar{r}^n = \bar{r}^m$ ならば $n \equiv m \pmod l$ となることを示せ．
\vspace{11cm}

\vfill\null
\end{multicols*}

\newpage

% ==========================================
% 1枚目 裏 (2ページ目: IMG_6468)
% ==========================================
\begin{multicols*}{2}
% 領域3
\noindent
\textbf{問4.}\quad $K$ を $\mathbb{C}$ の部分体とし，多項式環 $K[x]$ を考える．\\
(1)[6点] $f(x), g(x) \in K[x]$ に対し
\[ \deg(f(x) + g(x)) \le \max\{\deg f(x), \deg g(x)\} \]
であり，さらに $\deg f(x) \ne \deg g(x)$ なら等号が成立することを示せ．
\vspace{10cm}

\noindent
(2)[6点] $K[x]^{\times} = K \setminus \{0\}$ を示せ．
\vspace{10cm}

\vfill\null
\columnbreak

% 領域4
\noindent
(3)[6点] $f(x), g(x) \in K[x]$ で $g(x) \ne 0$ とする．このとき，
\[ f(x) = h(x)g(x) + r(x) \quad (\deg r(x) < \deg g(x)) \]
を満たす $h(x), r(x) \in K[x]$ が存在することを示せ．
\vspace{14cm}

\noindent
(4)[6点] 上で求めた $h(x), r(x)$ は一意的であることを示せ．

\vfill\null
\end{multicols*}

\newpage

% ==========================================
% 2枚目 表 (3ページ目: IMG_6469)
% ==========================================
\testheader

\begin{multicols*}{2}
% 領域5
\noindent
\textbf{問5.}[4点] $(\mathbb{Z}/13\mathbb{Z})^{\times}$ において $\bar{7}^k = \bar{a}$ となる $1 \le a \le 12$ を求めて下の表を埋めよ。\\[0.5em]
\begin{tabular}{|c|c|c|c|c|c|c|}
\hline
$k$ & 1 & 2 & 3 & 4 & 5 & 6 \\ \hline
$a$ &  &  &  &  &  & \\ \hline
$k$ & 7 & 8 & 9 & 10& 11& 12\\ \hline
$a$ &  &  &  &  &  &  \\ \hline
\end{tabular}
\vspace{4cm}

\noindent
\textbf{問6.}[6点] $20^{123}$ を $13$ で割った余りを求めよ．
\vspace{7.5cm}

\noindent
\textbf{問7.}[6点] $x^5 \equiv 2 \pmod{13}$ を満たす整数 $x$ を求めよ．
\vspace{8cm}

\vfill\null
\columnbreak

% 領域6
\noindent
\textbf{問8.}\quad $a_1 = 154, a_2 = 231, a_3 = 210$ に対して以下の問いに答えよ．\\
(1)[3点] $d = \operatorname{GCD}(a_1, a_2, a_3)$ を求めよ．\\
(2)[6点] $d = k_1 a_1 + k_2 a_2 + k_3 a_3$ を満たす整数 $k_1, k_2, k_3$ を一組求めよ．

\vfill\null
\end{multicols*}

\newpage

% ==========================================
% 2枚目 裏 (4ページ目: IMG_6470)
% ==========================================
\begin{multicols*}{2}
% 領域7
\noindent
\textbf{問9.}\quad $\omega = \frac{-1 + \sqrt{-3}}{2}$ とおく．このとき $\mathbb{Z}[\omega] = \{a + b\omega \mid a, b \in \mathbb{Z}\}$ は $\mathbb{C}$ の部分環となることがいえる．また $\alpha \in \mathbb{C}$ に対し、$N(\alpha) = |\alpha|^2$ と定める．次の問いに答えよ．\\[0.5em]
(1)[4点] $\alpha = a + b\omega \in \mathbb{Z}[\omega]$ に対し $N(\alpha)$ を $a, b$ の式で表せ．
\vspace{6.5cm}

\noindent
(2)[6点] $N(\alpha) = 3$ となる $\alpha \in \mathbb{Z}[\omega]$ を一つ求めよ．
\vspace{7.5cm}

\noindent
(3)[6点] 単元の集合 $\mathbb{Z}[\omega]^{\times}$ を求めよ．
\vspace{8.5cm}

\vfill\null
\columnbreak

% 領域8
\noindent
\textbf{問10.}[8点] 有理式 $\displaystyle P(x) = \frac{1}{(x+1)^2(x^2+1)}$ を有理数の範囲で部分分数展開せよ．
\vspace{24cm}

\vfill\null
\end{multicols*}

\end{document}